\documentclass[11pt]{article}

\usepackage{xparse,environ}
\usepackage{adjustbox}

% \newcommand{\answerDirectory}{solutions}

../../macros/macros.tex

% Local override, kept so this homework's PDF matches its original exactly:
% the shared macros (src/macros/macros.tex) now define \inputAnswer differently (which also calls \refreshHeader).
\renewcommand{\inputAnswer}[1]{\input{\answerDirectory/#1}}


\inputAnswer{problem_0_answer}

\doHwHeader{Homework \BLUE{2}}

%%%%% EDIT THE FILES NAMED problem_X_answer.tex.  DON'T EDIT THIS FILE

%%%%%%%%%%%%%%%%%%%%%%%%%%%%%%%%%%%%%%%%%%%%%%%%%%%%%%%%%%%% 
%%%%%%%%%%%%%%%%%%%%%%%%%%%%%%%%%%%%%%%%%%%%%%%%%%%%%%%%%%%% 
%%%%%%%%%%%%%%%%%%%%%%%%%%%%%%%%%%%%%%%%%%%%%%%%%%%%%%%%%%%% 
%%%%%%%%%%%%%%%%%%%%%%%%%%%%%%%%%%%%%%%%%%%%%%%%%%%%%%%%%%%% 

\begin{document}


%%%%% EDIT THE FILES NAMED problem_X_answer.tex.  DON'T EDIT THIS FILE 

\paragraph{General instructions.} First read \BLUE{Lecture Note 4},
then answer the problems given in the rest of this document.
Edit this LaTeX template to add your answers, as described in the file named \texttt{00\_README.txt},
then compile it to produce \BLUE{\texttt{homework\_2.pdf}}\ifGradescope{} for submission to gradescope\fi.

\ifGradescope
\BLUE{
This time when you submit to gradescope you don't need to identify the pages
for each problem.  We're using gradescope's ``fixed-length'' grading mode.
To make it work, your answer for each problem (and each part of Problem 3), should use at most 1 page.
Your PDF should have 10 pages.  Each answer should be in its designated location
(per the comment with each answer) in the PDF.
}
\fi

\paragraph{Don't know?  Say so.}
If you don't know the answer to a problem or any part of the problem,
we encourage you to simply write \emph{``I don't know''} instead of giving an answer.
(If you write ``I don't know'' you can also add any questions you might have.)

\paragraph{Gentle reminder: homework is for learning.}
As noted in Lecture Note 1, engaging fully with the homeworks is the best way to learn the material.
\emph{The primary role of homeworks is to help you learn the material,
  and the primary role of feedback on homeworks is to help you improve your ideas.}

Please try to explain your ideas as clearly as you can.
The clearer your explanations are, the more useful the feedback on them can be.

\paragraph{Collaboration and citation policy.}
If you are doing this homework as part of a course, follow the course's collaboration and citation policy.
In any case, at the end of each answer, list the people you collaborated with,
and cite every source that your answer uses ideas from
(other than the lecture notes and the sources listed in the lecture notes).
If there are none, write ``none''.

%%%%% EDIT THE FILES NAMED problem_X_answer.tex.  DON'T EDIT THIS FILE 

\vfill
\noindent{\footnotesize\copyrightNotice}

%%% Local Variables:
%%% mode: latex
%%% TeX-master: "homework_2"
%%% End:


%%%%%%%%%%%%%%%%%%%%%%%%%%%%%%%%%%%%%%%%%%%%%%%%%%%%%%%%%%%% 
%%%%%%%%%%%%%%%%%%%%%%%%%%%%%%%%%%%%%%%%%%%%%%%%%%%%%%%%%%%% 
%%%%%%%%%%%%%%%%%%%%%%%%%%%%%%%%%%%%%%%%%%%%%%%%%%%%%%%%%%%% 
%%%%%%%%%%%%%%%%%%%%%%%%%%%%%%%%%%%%%%%%%%%%%%%%%%%%%%%%%%%% 

%%%%% EDIT THE FILES NAMED problem_X_answer.tex.  DON'T EDIT THIS FILE 

%%%%%%%%%%%%%%%% PROBLEM 1  %%%%%%%%%%%%%%%%%

\newpage 

{\small
\problem \emph{(Section 1.4)}
Reorder this sequence so each function is big-O of the following function:
\newcommand{\sep}{~~~~}

\[
  n^2
  \sep
  6^{\log_2 n}
  \sep 
  n!
  \sep 
  11^n
  \sep 
  \log_2 (n!)  
  \sep 
  2^{3n}
  \sep 
  10 n^{3/2}
  \sep 
  100 n
  \sep 
  3^{\log_4 n}
  \sep 
  \log_2^4 n
\]
No justification is required.  Give the reordered sequence below.  For grading, order the functions vertically, one on each line.
}

\setlength{\ITEMHT}{0.45in}
\injectEnumerateBlue 

\inputAnswer{problem_1_answer}

\restoreEnumerateBlue

%%%%%%%%%%%%%%%%%%%%%%%%%%%%%%%%%%%%%%%%%%%%%%%%%%%%%%%%%%%% 
%%%%%%%%%%%%%%%%%%%%%%%%%%%%%%%%%%%%%%%%%%%%%%%%%%%%%%%%%%%% 
%%%%%%%%%%%%%%%%%%%%%%%%%%%%%%%%%%%%%%%%%%%%%%%%%%%%%%%%%%%% 
%%%%%%%%%%%%%%%%%%%%%%%%%%%%%%%%%%%%%%%%%%%%%%%%%%%%%%%%%%%% 

%%%%% EDIT THE FILES NAMED problem_X_answer.tex.  DON'T EDIT THIS FILE 

%%%%%%%%%%%%%%%% PROBLEM 2  %%%%%%%%%%%%%%%%%

\newpage 
{\small
\problem \emph{(Section 4.1)}
Use the so-called naive upper and lower bounds, and the rule for geometric sums,
to determine the big-$\Theta$ value of each of the following sums.
You don't have to show your work,
but you can if you want feedback on it
(and it may help you get partial credit if your answer is wrong). 

\begin{enumerate}[label=(\alph*)]

\item $\displaystyle\sum_{i=1}^{\log_2 n} i(\log_3 i)^2$

\item $\displaystyle\sum_{i=1}^{n} i^2/3^i$

\item $\displaystyle\sum_{i=1}^{2^n} (\log_2 i)^2$
  
\item $\displaystyle\sum_{i=1}^n 1.2^i/i$

\end{enumerate}
}

\setlength{\ITEMHT}{0.95in}
{
\injectEnumerate 

\inputAnswer{problem_2_answer}
}
% \restoreEnumerate 

%%%%%%%%%%%%%%%%%%%%%%%%%%%%%%%%%%%%%%%%%%%%%%%%%%%%%%%%%%%% 
%%%%%%%%%%%%%%%%%%%%%%%%%%%%%%%%%%%%%%%%%%%%%%%%%%%%%%%%%%%% 
%%%%%%%%%%%%%%%%%%%%%%%%%%%%%%%%%%%%%%%%%%%%%%%%%%%%%%%%%%%% 
%%%%%%%%%%%%%%%%%%%%%%%%%%%%%%%%%%%%%%%%%%%%%%%%%%%%%%%%%%%% 

%%%%% EDIT THE FILES NAMED problem_X_answer.tex.  DON'T EDIT THIS FILE 

%%%%%%%%%%%%%%%% PROBLEM 3  %%%%%%%%%%%%%%%%%

\newpage


\newcommand{\ALG}[4]{
  \vbox{                        % keep on one page
    \begin{steps}
    \item[] def \textsf{Print{#1}s}$(n)$:  \comment{Assume $n$ is a power of $#4$.}
      \step for $i = 1$ to $#2$: print(``{#1}'') \comment{\ldots print $#2$ #1's}
      \step if $n > 1$:
      \step~~~for $j=1$ to $#3$: \textsf{Print{#1}s}$(\floor{n/#4})$
      \comment{\ldots make $#3$ recursive calls, each with input $\floor{n/#4}$}
    \end{steps}
  }
}

\newpage

\problem\label{problem: recursion}\emph{(Recursion-tree method)}
Repeat Exercise 4.2 for each of the algorithms below.
Give your answers in the same format as the solution in Lecture Note 4 for that exercise.
To make it easier, ignore floors in your calculations.
That is, for (a) and (c) assume $n$ is a power of 3.
For (b) assume $n$ is a power of 4. 

\begin{enumerate}[label=(\alph*)]
\item \ALG{A}{3n}{4}{3}
  
\item \ALG{B}{n^2}{16}{4}

\item \ALG{C}{5n}{2}{3}

\end{enumerate}

\continued 

\setlength{\ITEMHT}{1.3in}
{
\injectEnumerate 

\inputAnswer{problem_3a_answer}

\continued 

\inputAnswer{problem_3b_answer}

\continued 

\inputAnswer{problem_3c_answer}
}
% \restoreEnumerate

%%%%%%%%%%%%%%%%%%%%%%%%%%%%%%%%%%%%%%%%%%%%%%%%%%%%%%%%%%%% 
%%%%%%%%%%%%%%%%%%%%%%%%%%%%%%%%%%%%%%%%%%%%%%%%%%%%%%%%%%%% 
%%%%%%%%%%%%%%%%%%%%%%%%%%%%%%%%%%%%%%%%%%%%%%%%%%%%%%%%%%%% 
%%%%%%%%%%%%%%%%%%%%%%%%%%%%%%%%%%%%%%%%%%%%%%%%%%%%%%%%%%%% 

%%%%% EDIT THE FILES NAMED problem_X_answer.tex.  DON'T EDIT THIS FILE 

%%%%%%%%%%%%%%%% PROBLEM 4  %%%%%%%%%%%%%%%%%

\newpage

{\small
\problem\label{problem: mergesort}\emph{(recursion trees)}
Consider a variant of \textsf{mergesort} that, given an array of size $N$,
recurses as usual unless the subproblem to be solved is of size $\sqrt N$ or less,
in which case it solves the subproblem without recursing, instead using insertion sort
(a quadratic-time algorithm) on the subproblem.
Use a recursion-tree analysis to determine
the asymptotic worst-case running time of this algorithm on arrays of size $N$.
Show your work as described for Problem~\ref{problem: recursion}.

You may assume $N$ is an even power of two.  That is $N=2^{2k}$ for some integer $k$,
so $\sqrt N = 2^k$.
}

\setlength{\ITEMHT}{1.15in}
{
\injectEnumerate 

\inputAnswer{problem_4_answer}
}
% \restoreEnumerate

%%%%%%%%%%%%%%%%%%%%%%%%%%%%%%%%%%%%%%%%%%%%%%%%%%%%%%%%%%%% 
%%%%%%%%%%%%%%%%%%%%%%%%%%%%%%%%%%%%%%%%%%%%%%%%%%%%%%%%%%%% 
%%%%%%%%%%%%%%%%%%%%%%%%%%%%%%%%%%%%%%%%%%%%%%%%%%%%%%%%%%%% 
%%%%%%%%%%%%%%%%%%%%%%%%%%%%%%%%%%%%%%%%%%%%%%%%%%%%%%%%%%%% 

%%%%% EDIT THE FILES NAMED problem_X_answer.tex.  DON'T EDIT THIS FILE 

%%%%%%%%%%%%%%%% PROBLEM 5 %%%%%%%%%%%%%%%%%

\newpage

{\small

\problem\emph{(Exercise 3.13)}\label{problem: busy hospitals}
\textbf{NOTE: Do this problem or Problem~\ref{problem: lonely doctor}, but not both.}

Prove or disprove the following theorem:

\begin{theorem}
  In every execution of the hospitals-propose \prob{Stable Matching} algorithm,
  there is at most one hospital that makes offers to every doctor.
\end{theorem}

(Note: we mean the algorithm described in the lecture notes
  --- in each iteration, one hospital makes an offer to one doctor.)
If the theorem is false, describe a counter-example:
a particular instance, and an execution of the algorithm on that instance,
% [review 2026-09-27] fix: "make offers to every hospital" -> "every doctor" (hospitals make offers to doctors; matches LN3 Exercise 3.13)
in which at least two hospitals make offers to every doctor.
If the theorem is true, give a long-form proof.
In the proof, you may reuse without proof any observations, lemmas, or theorems from the lecture notes,
but cite the observation, lemma, or theorem by number when you do, to let the reader know.
}
%%%%%%%%%%%%%%%%%%%%%%%%%%%%%%%%%%%%%%%%%%%%%%%%%%%%%%%%%%%% 
%%%%%%%%%%%%%%%%%%%%%%%%%%%%%%%%%%%%%%%%%%%%%%%%%%%%%%%%%%%% 

%%%%% EDIT THE FILES NAMED problem_X_answer.tex.  DON'T EDIT THIS FILE 

\inputAnswer{problem_5_answer}


%%%%%%%%%%%%%%%%%%%%%%%%%%%%%%%%%%%%%%%%%%%%%%%%%%%%%%%%%%%% 
%%%%%%%%%%%%%%%%%%%%%%%%%%%%%%%%%%%%%%%%%%%%%%%%%%%%%%%%%%%% 
%%%%%%%%%%%%%%%%%%%%%%%%%%%%%%%%%%%%%%%%%%%%%%%%%%%%%%%%%%%% 
%%%%%%%%%%%%%%%%%%%%%%%%%%%%%%%%%%%%%%%%%%%%%%%%%%%%%%%%%%%% 

%%%%% EDIT THE FILES NAMED problem_X_answer.tex.  DON'T EDIT THIS FILE 

%%%%%%%%%%%%%%%% PROBLEM 6 %%%%%%%%%%%%%%%%%

\newpage
{\small
\problem\emph{(Exercise 3.11)}\label{problem: lonely doctor}
\textbf{NOTE: Do this problem or Problem~\ref{problem: busy hospitals}, but not both.}

Prove or disprove the following theorem:

\begin{theorem}
  In any execution of the hospitals-propose \prob{Stable Matching} algorithm
  there is at least one doctor that receives no offer before the final iteration.
\end{theorem}

Note: we mean the algorithm described in the lecture notes
  --- in each iteration, one hospital makes an offer to one doctor.)
If the theorem is false, describe a counter-example:
a particular instance, and an execution of the algorithm on that instance,
in which every doctor receives an offer before the final iteration.
If the theorem is true, give a long-form proof.
In the proof, you may reuse without proof any observations, lemmas, or theorems from the lecture notes,
but cite the observation, lemma, or theorem by number when you do, to let the reader know.
}
%%%%%%%%%%%%%%%%%%%%%%%%%%%%%%%%%%%%%%%%%%%%%%%%%%%%%%%%%%%% 
%%%%%%%%%%%%%%%%%%%%%%%%%%%%%%%%%%%%%%%%%%%%%%%%%%%%%%%%%%%% 

%%%%% EDIT THE FILES NAMED problem_X_answer.tex.  DON'T EDIT THIS FILE 

\inputAnswer{problem_6_answer}


\end{document}

%%% Local Variables:
%%% mode: latex
%%% TeX-master: t
%%% End:
